\documentclass[12pt,a4paper]{article}

% Packages
\usepackage[margin=1in]{geometry}
\usepackage{graphicx}
\usepackage{hyperref}
\usepackage{listings}
\usepackage{xcolor}
\usepackage{booktabs}
\usepackage{amsmath}
\usepackage{fancyhdr}
\usepackage{float}
\usepackage{longtable}
\usepackage{array}
\usepackage{parskip}
\usepackage{caption}
\usepackage{adjustbox}

% Code listing style
\definecolor{codeblue}{RGB}{30,100,200}
\definecolor{codegray}{RGB}{100,100,100}
\definecolor{codegreen}{RGB}{0,130,60}
\definecolor{codebg}{RGB}{248,248,248}

\lstdefinestyle{apistyle}{
    backgroundcolor=\color{codebg},
    commentstyle=\color{codegreen},
    keywordstyle=\color{codeblue}\bfseries,
    numberstyle=\tiny\color{codegray},
    stringstyle=\color{orange},
    basicstyle=\ttfamily\small,
    breaklines=true,
    captionpos=b,
    keepspaces=true,
    numbers=left,
    numbersep=5pt,
    frame=single,
    rulecolor=\color{gray!30}
}
\lstset{style=apistyle}

% Header/Footer
\pagestyle{fancy}
\fancyhf{}
\rhead{Village Pond Planner}
\lhead{CSD Assignment Report}
\rfoot{Page \thepage}
\lfoot{Roll No: stu84}

\hypersetup{colorlinks=true, linkcolor=blue, urlcolor=blue}

\begin{document}

% ─── Title Page ───────────────────────────────────────────────────────────────
\begin{titlepage}
    \centering
    \vspace*{2cm}
    {\huge\bfseries Village Pond Planner\\[0.5em]}
    {\Large Contour-Driven Hydrological Analysis API\\[2em]}
    \rule{\linewidth}{0.5pt}\\[1em]
    {\large\textbf{CSD Assignment Report}\\[0.5em]}
    {\large Phase 1: Catchment Delineation \& Pond Siting API}
    \rule{\linewidth}{0.5pt}\\[2em]
    \begin{tabular}{rl}
        \textbf{Student:}    & Rithika Sannidhi \\
        \textbf{Roll No:}    & stu84 \\
        \textbf{GitHub:}     & \url{https://github.com/rithika1503/village-pond-planner} \\
        \textbf{API URL:}    & \texttt{http://10.1.75.53:4293} \\
        \textbf{Swagger UI:} & \texttt{http://10.1.75.53:4293/docs} \\
        \textbf{Date:}       & \today \\
    \end{tabular}
    \vfill
    {\small Department of Computer Science}
\end{titlepage}

\tableofcontents
\newpage

% ─── 1. Introduction ──────────────────────────────────────────────────────────
\section{Introduction}

This report documents the design, implementation, and deployment of the \textbf{Village Pond Planner} --- a REST API that accepts an uploaded contour map in KML format, derives the underlying terrain hydrology using a grid-and-graph based algorithm, identifies optimal pond locations that are safely away from rivers, and returns a structured JSON response with catchment geometry, pond sizing, and suitability scores.

The system is \textbf{fully general}: no hard-coded coordinates or pre-computed results are embedded. All outputs are derived algorithmically from the uploaded contour file, making the system extensible to any region worldwide.

% ─── 2. Algorithm Approach ────────────────────────────────────────────────────
\section{Algorithm Approach}

\subsection{Overview: Grid, Graph, and Flow Routing}

The core idea is to convert a sparse set of contour lines into a \textbf{regular grid} (a Digital Elevation Model, or DEM), then model that grid as a \textbf{directed graph} where every cell points to its downhill neighbour. Classical graph algorithms --- a priority queue, a topological sort, and Breadth-First Search (BFS) --- are then applied on this graph to simulate how water flows and accumulates across the terrain.

\subsection{Step 1: Building the Grid (DEM Interpolation)}

The uploaded KML file contains hundreds of contour lines, each a polyline with a known elevation. These are sparse, irregular data points. Using \texttt{scipy.interpolate.griddata} with cubic interpolation, all these contour points are interpolated onto a uniform \textbf{150$\times$150 grid} of latitude-longitude cells. Each cell in this grid stores one elevation value. This grid is the DEM, and it forms the foundation of every subsequent computation.

\subsection{Step 2: Sink Filling (Priority Queue)}

A raw DEM often has \textbf{sinks} --- cells lower than all their neighbours, which would cause water to pool and stop flowing instead of draining to the boundary. To fix this, the Wang \& Liu (2006) algorithm is applied:

\begin{enumerate}
    \item All boundary cells are pushed into a \textbf{min-heap priority queue} with their true elevations as priority keys.
    \item The lowest cell is popped; its unvisited neighbours are given a \textit{spill elevation} equal to the maximum of their own elevation and the current cell's spill elevation, and are pushed into the heap.
    \item This ensures every interior cell has a guaranteed downhill path to the boundary --- sinks are filled by raising them to the level at which they would naturally overflow.
\end{enumerate}

After this step, the DEM is hydrologically correct: every cell drains continuously to the edge.

\subsection{Step 3: D8 Flow Direction (Graph Edge Assignment)}

Each DEM cell is now treated as a \textbf{node in a directed graph}. For every interior cell, the steepest downhill neighbour among its 8 surrounding cells (cardinal + diagonal) is found, and a directed edge is drawn from the current cell to that neighbour:

\begin{equation}
    \text{slope}_{d} = \frac{\Delta z}{\Delta d}, \quad
    \Delta d = \begin{cases} 1 & \text{cardinal direction} \\ \sqrt{2} & \text{diagonal direction} \end{cases}
\end{equation}

\noindent where $\Delta z$ is the elevation difference and $\Delta d$ is the cell-to-cell distance. The result is a \textbf{flow direction graph}: each cell has exactly one outgoing edge pointing to where water would flow next.

\subsection{Step 4: Flow Accumulation (Topological Sort)}

With the directed flow graph constructed, the \textbf{upstream contributing area} of every cell is computed using a topological sort. Cells are processed in reverse flow order (from ridges to valleys): each cell's accumulation count is the number of cells that drain through it. Cells at river channels naturally accumulate the most upstream area, while ridge cells have zero accumulation.

\subsection{Step 5: River Exclusion Using OpenStreetMap}

A critical requirement is to \textbf{never place a pond inside a river}. Two complementary masking strategies are applied:

\begin{itemize}
    \item \textbf{Mathematical mask:} Any grid cell whose flow accumulation exceeds the 85th percentile is classified as an active stream channel and permanently excluded from pond candidates. These are the cells that drain the most water and correspond to the main river thalwegs.

    \item \textbf{OpenStreetMap (OSM) physical mask:} The official OpenStreetMap REST API (\texttt{api.openstreetmap.org/api/0.6/map}) is queried for all physical waterway features (rivers, streams, canals) within the bounding box of the input map. The XML response is parsed using Python's \texttt{xml.etree.ElementTree} to extract each waterway as a sequence of latitude/longitude nodes. Each waterway polyline is then rasterised onto the DEM grid: any grid cell within \textbf{220 metres} of a physical waterway is permanently excluded. This ensures that even wide river valleys and meandering channels are fully masked using real-world geographic data.
\end{itemize}

Together, these two masks ensure that no recommended pond location can fall inside a river channel, whether detected mathematically or physically from OpenStreetMap.

\subsection{Step 6: Candidate Identification (BFS Region Growing)}

After exclusion, the remaining grid cells are filtered to find those that are geologically suitable:
\begin{itemize}
    \item Slope $< 5°$ (gentle terrain, practical for earthwork)
    \item Elevation $\leq$ 40th percentile (valley position, naturally collects runoff)
    \item Not masked by either river exclusion criterion above
\end{itemize}

Adjacent qualifying cells are grouped into contiguous \textbf{candidate regions} using a \textbf{4-connectivity BFS flood-fill}: starting from any qualifying cell, all directly-adjacent qualifying cells are discovered and merged into one region. The representative cell chosen for each region is the one with the highest flow accumulation --- the natural hydrological convergence point of that region where water pools most reliably.

A \textbf{500\,m spatial-diversity filter} ensures no two final candidates are closer than 500 m to each other, producing geographically spread options for the planner.

\subsection{Step 7: Catchment Delineation (BFS Upstream Trace)}

For each candidate outlet, the system performs an \textbf{upstream BFS} on the flow direction graph: starting at the outlet cell, it traverses all cells whose flow paths eventually lead to that outlet. The set of all such cells constitutes the \textbf{catchment area} --- every drop of rainfall that falls inside this boundary will flow to the pond. The catchment polygon is constructed by taking the union of all contributing grid cells.

\subsection{Step 8: Pond Sizing (Rational Method)}

Annual runoff and pond sizing are estimated using the Rational Method:

\begin{align}
    V_{\text{runoff}} &= P \times A \times C \\[4pt]
    A_{\text{pond}} &= \frac{V_{\text{runoff}}}{d_{\text{pond}}}
\end{align}

\noindent where:
\begin{itemize}
    \item $V_{\text{runoff}}$ = estimated annual runoff volume (m$^3$/yr)
    \item $P$ = annual rainfall in metres per year, fetched live from the Open-Meteo ERA5 climate archive
    \item $A$ = catchment area (m$^2$), computed from the upstream BFS in Step 7
    \item $C$ = dimensionless runoff coefficient representing land cover type:\\
    \hspace*{1em} $C = 0.30$ (vegetated), $C = 0.55$ (agricultural), $C = 0.75$ (built-up)
    \item $A_{\text{pond}}$ = required pond surface area (m$^2$)
    \item $d_{\text{pond}}$ = desired pond depth in metres (default 3\,m)
\end{itemize}

\subsection{Step 9: Suitability Scoring}

Each candidate is ranked using a weighted composite score:

\begin{equation}
    S = w_T \cdot T + w_C \cdot C + w_L \cdot L
\end{equation}

\noindent where:
\begin{itemize}
    \item $S$ = overall suitability score (0 to 1)
    \item $T$ = terrain score: normalised combination of slope and relative elevation
    \item $C$ = catchment adequacy score: based on catchment area relative to a target storage volume
    \item $L$ = land availability score: reflects suitability of land cover type
    \item $w_T = 0.40$, $w_C = 0.40$, $w_L = 0.20$ (weights sum to 1.0)
\end{itemize}

% ─── 3. API Documentation ─────────────────────────────────────────────────────
\section{API Route Documentation}

\subsection{Endpoints}

\begin{center}
\begin{tabular}{p{1.2cm}p{3.8cm}p{8cm}}
    \toprule
    \textbf{Method} & \textbf{Path} & \textbf{Description} \\
    \midrule
    POST & \texttt{/analyzeContour} & Primary endpoint: accepts KML/KMZ, returns full analysis JSON \\
    POST & \texttt{/findCatchment}  & Alias for \texttt{/analyzeContour} \\
    GET  & \texttt{/docs}           & Interactive Swagger UI \\
    GET  & \texttt{/health}         & Health check \\
    \bottomrule
\end{tabular}
\end{center}

\subsection{Request Parameters}

\begin{center}
\begin{tabular}{p{3.5cm}p{1.5cm}p{1.5cm}p{6cm}}
    \toprule
    \textbf{Field} & \textbf{Type} & \textbf{Required} & \textbf{Description} \\
    \midrule
    \texttt{file} & File & Yes & KML or KMZ contour map \\
    \texttt{land\_cover} & string & No & \texttt{vegetated / agricultural /} \texttt{built\_up}. Default: \texttt{agricultural} \\
    \texttt{desired\_depth\_m} & float & No & Target pond depth in metres. Default: 3.0 \\
    \texttt{annual\_rainfall\_mm} & float & No & Override rainfall. Fetched from Open-Meteo if omitted \\
    \bottomrule
\end{tabular}
\end{center}

\subsection{cURL Example}

\begin{lstlisting}[language=bash, caption={Full API call with the sample contour map}]
curl -X POST "http://10.1.75.53:4293/analyzeContour" \
  -F "file=@contours_1m.kml" \
  -F "land_cover=agricultural" \
  -F "desired_depth_m=3.0" \
  | python3 -m json.tool
\end{lstlisting}

% ─── 4. Demonstration ─────────────────────────────────────────────────────────
\section{Demonstration with Sample Contour Map}

\subsection{Input}
\begin{itemize}
    \item File: \texttt{contours\_1m.kml} (2710 contour lines)
    \item Region: Bilaspur, Chhattisgarh, India
    \item Bounding Box: $21.2398^\circ$--$21.2636^\circ$N, $81.2814^\circ$--$81.3126^\circ$E
    \item Elevation range: 267\,m to 298\,m
    \item OSM waterways fetched and excluded: 5 features, 6986 grid cells masked
\end{itemize}

\begin{figure}[H]
    \centering
    \fbox{\rule{0pt}{6cm}\rule{12cm}{0pt}}
    \caption{Filled DEM visualisation showing elevation gradient (blue = low/valley, yellow/red = high/ridge) with top candidate locations marked.
    \textit{[Insert \texttt{filled\_dem.png} from \texttt{python scripts/debug\_flow.py}]}}
    \label{fig:filled_dem}
\end{figure}

\begin{figure}[H]
    \centering
    \fbox{\rule{0pt}{5cm}\rule{12cm}{0pt}}
    \caption{Flow accumulation map. Bright cells = major river channels (high upstream area). Candidate sites are placed on dark cells (valley margins) safely away from the bright river centre-lines.
    \textit{[Insert \texttt{flow\_accumulation.png} from \texttt{python scripts/debug\_flow.py}]}}
    \label{fig:flow_accum}
\end{figure}

\subsection{Top 10 Candidate Pond Sites}

\begin{table}[H]
\centering
\caption{Top 10 Pond Candidate Sites derived from \texttt{contours\_1m.kml}}
\label{tab:top10}
\small
\begin{adjustbox}{max width=\textwidth}
\begin{tabular}{clllcccc}
    \toprule
    \textbf{Rank} & \textbf{Latitude} & \textbf{Longitude} & \textbf{Elev (m)} & \textbf{Slope ($^\circ$)} & \textbf{Flow Acc.} & \textbf{Terrain Score} & \textbf{Catchment (ha)} \\
    \midrule
    1  & 21.2524 & 81.2814 & 280.0 & 0.53 & 14 & 0.95 & 5.08   \\
    2  & 21.2413 & 81.2897 & 279.0 & 0.34 & 13 & 0.95 & 124.59 \\
    3  & 21.2616 & 81.2967 & 282.0 & 0.92 & 14 & 0.90 & 8.88   \\
    4  & 21.2558 & 81.2897 & 279.6 & 1.25 & 13 & 0.87 & 1.90   \\
    5  & 21.2602 & 81.2884 & 278.0 & 1.99 & 14 & 0.84 & 33.90  \\
    6  & 21.2413 & 81.2827 & 278.5 & 2.26 & 14 & 0.82 & 4.73   \\
    7  & 21.2447 & 81.2999 & 280.8 & 0.78 &  9 & 0.78 & 69.54  \\
    8  & 21.2461 & 81.2865 & 277.9 & 2.77 & 14 & 0.78 & 16.02  \\
    9  & 21.2558 & 81.2954 & 281.0 & 0.53 &  8 & 0.77 & 15.82  \\
    10 & 21.2573 & 81.2827 & 279.0 & 0.93 &  8 & 0.76 & 4.17   \\
    \bottomrule
\end{tabular}
\end{adjustbox}
\end{table}

\begin{figure}[H]
    \centering
    \fbox{\rule{0pt}{6cm}\rule{12cm}{0pt}}
    \caption{Top 10 candidate sites with catchment polygons on satellite basemap (geojson.io).
    \textit{[Insert screenshot of \texttt{top10\_visualization.geojson} opened on geojson.io]}}
    \label{fig:top10}
\end{figure}

\subsection{Why These Sites Are Best}

\subsubsection{Site \#1 --- Primary Recommendation}
\textbf{Location:} $21.2524^\circ$N, $81.2814^\circ$E \quad\textbf{Catchment: 5.08\,ha}

Highest terrain score (0.95) due to near-flat slope ($0.53^\circ$) and low valley position. The BFS upstream trace confirms this cell is the hydrological convergence point for a 5.08\,ha sub-watershed --- meaning almost all rainfall in that area flows directly here. Minimal excavation required.

\subsubsection{Site \#2 --- Largest Catchment}
\textbf{Location:} $21.2413^\circ$N, $81.2897^\circ$E \quad\textbf{Catchment: 124.59\,ha}

The gentlest slope of all candidates ($0.34^\circ$) and the largest catchment (124.59\,ha) make this site the highest water-yield option. Suitable for a large village reservoir that must serve irrigation. The large catchment is a result of a major tributary confluence detected by the flow accumulation graph.

\subsubsection{Site \#3 --- Northern Valley}
\textbf{Location:} $21.2616^\circ$N, $81.2967^\circ$E \quad\textbf{Catchment: 8.88\,ha}

Located in the northern portion of the study area, this site provides a geographically independent option to Site \#1. Its moderate 8.88\,ha catchment is well-suited for a medium-scale community pond. Good terrain score (0.90) confirms reliable water collection.

\subsubsection{Sites \#4 to \#10 --- Distributed Alternatives}

Sites 4 through 10 represent viable alternatives spread across the study region, with catchments ranging from 1.90\,ha to 69.54\,ha. The 500\,m spatial-diversity filter guarantees that each of these sites serves a geographically distinct community or sub-watershed. Site \#7 (69.54\,ha catchment) is a notable secondary large-yield option in the eastern portion of the area.

\begin{figure}[H]
    \centering
    \fbox{\rule{0pt}{6cm}\rule{12cm}{0pt}}
    \caption{Top 3 candidate sites with their individual catchment polygons.
    \textit{[Insert screenshot of \texttt{top3\_visualization.geojson} opened on geojson.io]}}
    \label{fig:top3}
\end{figure}

\subsection{Summary Statistics}

\begin{center}
\begin{tabular}{ll}
    \toprule
    \textbf{Metric} & \textbf{Value} \\
    \midrule
    Best Pond Location & $21.2524^\circ$N, $81.2814^\circ$E \\
    Best Site Catchment & 5.08\,ha \\
    Largest Catchment (Site \#2) & 124.59\,ha \\
    Annual Rainfall (ERA5) & 1333\,mm/yr \\
    Runoff Volume (Best Site) & 30,375\,m$^3$/yr \\
    Pond Surface Area & 10,125\,m$^2$ \\
    Pond Depth & 3.0\,m \\
    Suitability Score & 0.73 / 1.0 \\
    Candidates Returned & 10 \\
    OSM Waterways Masked & 5 features, 6986 cells \\
    \bottomrule
\end{tabular}
\end{center}

% ─── 5. Extensibility ─────────────────────────────────────────────────────────
\section{Code Extensibility}

\begin{itemize}
    \item \textbf{New input format:} Only \texttt{kml\_parser.py} needs modification to return \texttt{(lats, lons, elevs)} arrays. All downstream modules are reused unchanged.
    \item \textbf{New scoring dimension:} Add a weight to \texttt{config.py} and update \texttt{scoring.py}. No router changes needed.
    \item \textbf{Raster DEM input:} Replace \texttt{contours\_to\_dem()} with a rasterio reader returning the same NumPy arrays.
    \item \textbf{Multi-village:} The API is fully stateless --- call it $N$ times with $N$ KML files for $N$ independent analyses.
    \item \textbf{All parameters tunable} in \texttt{backend/app/config.py}: grid resolution, slope threshold, stream percentile, OSM buffer, scoring weights.
\end{itemize}

% ─── 6. GitHub ────────────────────────────────────────────────────────────────
\section{GitHub Repository}

\begin{center}
    \url{https://github.com/rithika1503/village-pond-planner}
\end{center}

\begin{center}
\begin{tabular}{ll}
    \toprule
    \textbf{File} & \textbf{Role} \\
    \midrule
    \texttt{app/routers/contour.py} & API endpoint (\texttt{POST /analyzeContour}) \\
    \texttt{app/geospatial/hydrology\_engine.py} & Grid graph: sink-fill, D8, flow accumulation \\
    \texttt{app/geospatial/catchment.py} & BFS upstream catchment delineation \\
    \texttt{app/geospatial/terrain.py} & Candidate identification + OSM masking \\
    \texttt{app/geospatial/osm\_client.py} & Official OSM API XML client \\
    \texttt{app/config.py} & All tunable parameters \\
    \texttt{scripts/debug\_flow.py} & Generates DEM and flow PNG visualisations \\
    \texttt{scripts/generate\_top3\_geojson.py} & Generates GeoJSON for QGIS/geojson.io \\
    \bottomrule
\end{tabular}
\end{center}

% ─── 7. Deployment ────────────────────────────────────────────────────────────
\section{Deployment}

Deployed on \textbf{stu84\_sys1} (\texttt{10.1.75.53}). Per lab port rules (SSH port $22XX \Rightarrow$ App ports $3XX$--$7XX$), with SSH port 2293, the application runs on port \textbf{4293} (Application 2 slot).

\begin{lstlisting}[language=bash, caption={Deployment on stu84\_sys1 using tmux}]
# SSH in
ssh -p 2293 student@10.1.75.53

# Set up Python 3.12 environment (no compiled packages needed)
cd ~/village-pond-planner/village-pond-planner/backend
python3.12 -m venv venv --without-pip
source venv/bin/activate
curl -sS https://bootstrap.pypa.io/get-pip.py | python
pip install -r requirements_server.txt

# Start in tmux (keeps running after SSH disconnect)
tmux new-session -s pond
source venv/bin/activate
uvicorn app.main:app --host 0.0.0.0 --port 4293
# Press Ctrl+B, then D to detach
\end{lstlisting}

\begin{center}
\begin{tabular}{ll}
    \toprule
    \textbf{Resource} & \textbf{URL} \\
    \midrule
    Live API & \texttt{http://10.1.75.53:4293/analyzeContour} \\
    Swagger UI & \texttt{http://10.1.75.53:4293/docs} \\
    Health check & \texttt{http://10.1.75.53:4293/health} \\
    \bottomrule
\end{tabular}
\end{center}

\end{document}
